% !TeX program = lualatex
% Standalone notebook; compile twice: lualatex 35.tex
% Research batch 39 down to 25, 27 September 2026.
% Original section preserved verbatim except for marked insertion blocks.
\documentclass[11pt,a4paper]{article}
\usepackage[margin=24mm,headheight=14pt]{geometry}
\usepackage{fontspec}
\setmainfont{Latin Modern Roman}
\setsansfont{Latin Modern Sans}
\setmonofont{Latin Modern Mono}
\usepackage{amsmath,amssymb,mathtools}
\usepackage{unicode-math}
\setmathfont{Latin Modern Math}
\usepackage{microtype}
\usepackage{enumitem,xurl,needspace}
\usepackage[dvipsnames]{xcolor}
\usepackage{hyperref}
\hypersetup{unicode=true,colorlinks=true,linkcolor=MidnightBlue,urlcolor=MidnightBlue,
 pdftitle={Research notebook: Problem 35},
 pdfauthor={Open Problems Math research notebook},bookmarksnumbered=true}
\usepackage{bookmark,titlesec,fancyhdr}
\titleformat{\section}{\Large\bfseries\raggedright}{\thesection}{0.7em}{}
\pagestyle{fancy}\fancyhf{}
\fancyhead[L]{\small\sffamily Open Problems Math\enspace|\enspace Research notebook}
\fancyhead[R]{\small\sffamily Problem 35}
\fancyfoot[L]{\small\sffamily 27 September 2026}
\fancyfoot[R]{\small\sffamily \thepage}
\setlength{\parindent}{0pt}
\setlength{\parskip}{5pt plus 1pt}
\setlength{\emergencystretch}{3em}
\setcounter{secnumdepth}{1}
\setcounter{tocdepth}{1}
\setlist[itemize]{leftmargin=3em,itemsep=5pt,topsep=4pt,parsep=0pt}
\allowdisplaybreaks
\newcommand{\N}{\mathbb{N}}
\newcommand{\Z}{\mathbb{Z}}
\newcommand{\Q}{\mathbb{Q}}
\newcommand{\R}{\mathbb{R}}
\newcommand{\C}{\mathbb{C}}
\newcommand{\F}{\mathbb{F}}
\newcommand{\A}{\mathbb{A}}
\newcommand{\PP}{\mathbb P}
\newcommand{\E}{\mathbb{E}}
\newcommand{\eps}{\varepsilon}
\newcommand{\dd}{\,\mathrm d}
\newcommand{\sref}[2]{\hyperlink{source:#1:#2}{[S#2]}}
\newcommand{\eref}[2]{\hyperlink{extra:#1:#2}{[E#2]}}
\newif\ifshowcatalogue
\showcataloguetrue
\newcommand{\cataloguescope}[1]{\ifshowcatalogue
 \paragraph{Exact scope in the supplied catalogue.}
 \begin{quote}\small #1\end{quote}\fi}
\numberwithin{equation}{section}
\begin{document}
\setcounter{section}{34}
% BEGIN PRESERVED SECTION
% ================================================================
% problemId: problem.unconditional-separation-of-bqp-from-the-polynomial-hierarchy
% source releaseRank: 35; source releaseStatus: open
% exactTarget SHA-256: 20b53bd662f45ba7cc0c2ed1b04f13e37f9c896eef637e3c1ac8e1eea9862a8d
\clearpage
\section{\texorpdfstring{Unconditional Separation of BQP from the Polynomial Hierarchy}{Unconditional Separation of BQP from the Polynomial Hierarchy}}\label{35}
\subsection{Definitions and mathematical statement}
The polynomial hierarchy is $\mathsf{PH}=\bigcup_{k\ge0}\Sigma_k^{\mathsf P}$, where a language in $\Sigma_k^{\mathsf P}$ has a polynomial-time predicate preceded by $k$ alternating blocks of polynomially bounded existential and universal quantifiers, starting existentially. The target is
\[
 \exists L\in\mathsf{BQP}\quad L\notin\Sigma_k^{\mathsf P}\text{ for every fixed }k.
\]
This is a statement about ordinary uniform computation. Exhibiting an oracle relative to which the containment fails is a different statement.
\par\smallskip\textit{Formulation and concept references:} \sref{35}{1}.
\cataloguescope{Prove in the standard unrelativized model that bounded-error quantum polynomial time BQP is not contained in the polynomial hierarchy PH.}
\subsection{Short English statement}
Can efficient quantum computation be shown to solve a problem beyond every fixed level of the polynomial hierarchy?
% BEGIN RESEARCH 35
\subsection{Research attempt}

\paragraph{Current frontier.}
The established Raz--Tal result is an oracle separation
$\mathsf{BQP}^O\not\subseteq\mathsf{PH}^O$, obtained from a distributional
query problem; deleting the superscripts is not a valid inference.
\eref{35}{1}. Buzet--Chailloux's August 2026 version implements
2-Forrelation with restricted IQP resources and strengthens an oracle
separation, not the present unrelativized statement. \sref{35}{3},
\eref{35}{2}. Miloschewsky--Podder--Rudolph's September 2026 manuscript
constructs an oracle where quantum and deterministic language classes agree
while their promise versions do not. Its abstract makes particularly clear
why an oracle promise must not silently become a total language.
\sref{35}{4}, \eref{35}{3}. The recent statements were checked; their full
proofs were not independently audited. Aaronson's original formulation
remains the target here. \sref{35}{1}.

\paragraph{Attack route.}
A plausible route is to replace the random oracle functions in Forrelation
by uniformly evaluable, succinct functions, then prove that every fixed level
of the polynomial hierarchy fails on the resulting total decision problem.
This requires both a total gap convention and a lower bound that survives
access to the functions' descriptions. Query hardness by itself supplies
neither. I test two natural explicit replacements, Fourier-sparse functions
and quadratic phases, before attempting a more ambitious derandomization.
The calculations below are deductions in this notebook, not new claims about
the published oracle construction.

\paragraph{Attempt.}
Set $N=2^m$, take $f,g:\F_2^m\to\{-1,1\}$, and use the normalization
\begin{equation}\label{r35:phi}
 \Phi(f,g)=N^{-3/2}\sum_{x,y\in\F_2^m}
             f(x)(-1)^{x\cdot y}g(y).
\end{equation}
This is the inner product of two normalized states after a normalized
Hadamard transform, so $|\Phi|\le1$. With reversible evaluation circuits for
$f,g$, a Hadamard test estimates this real amplitude to fixed additive
accuracy with polynomially many gates in the description lengths. This
explains the attraction of a constant-gap promise; it does not define
bounded-error answers for inputs arbitrarily close to a threshold. The
Forrelation setup and its query interpretation are discussed in
\sref{35}{1} and \eref{35}{1}.

\textbf{Lemma 35.1 (a sparsity obstruction).}
Define $\widehat f(y)=N^{-1}\sum_x f(x)(-1)^{x\cdot y}$. If at most $k$
Fourier coefficients are nonzero, then
\begin{equation}\label{r35:sparse}
                    |\Phi(f,g)|\le\sqrt{k/N}.
\end{equation}
\emph{Proof.} Orthogonality gives $\sum_y\widehat f(y)^2=1$, and
$\Phi=N^{-1/2}\sum_y\widehat f(y)g(y)$. Cauchy--Schwarz on its support
proves the assertion. Thus $\Phi\ge c>0$ requires $k\ge c^2N$.
A family with polynomially many Fourier coefficients in $m$ cannot produce
constant positive Forrelation in this normalization. This says nothing
against succinct functions with exponentially large Fourier support.

\medskip
\textbf{Lemma 35.2 (quadratic phases admit exact classical evaluation).}
Suppose $f(x)=(-1)^{q(x)}$ and $g(y)=(-1)^{r(y)}$, with explicitly given
multilinear quadratic polynomials over $\F_2$. Then $\Phi(f,g)$ can be
computed exactly in time polynomial in $m$, including its sign.

\emph{Proof.} The numerator of \eqref{r35:phi} is
$S(Q)=\sum_{z\in\F_2^{2m}}(-1)^{Q(z)}$, where
$Q(x,y)=q(x)+x\cdot y+r(y)$ is quadratic. Here is an elimination algorithm
that does not enumerate the $2^{2m}$ points. If a cross term $uv$ occurs,
write, with $z$ denoting the other variables,
\[
 Q(u,v,z)=uv+uA(z)+vB(z)+C(z),
 \qquad \deg A,\deg B\le1,\quad\deg C\le2.
\]
For fixed $z$, direct summation gives
\begin{equation}\label{r35:eliminate}
 \sum_{u,v\in\F_2}(-1)^{uv+uA+vB+C}=2(-1)^{C+AB}.
\end{equation}
Replace $Q$ by $C+AB$, reducing $z_i^2=z_i$, and record a factor two.
The degree remains at most two. If no cross term remains, an affine
polynomial has sum zero when any linear coefficient is nonzero, and
otherwise sum $(-1)^{Q(0)}2^{\#z}$. At most $m$ eliminations occur;
coefficient bookkeeping uses polynomially many bit operations. The answer
is an integer divided by $2^{3m/2}$, representable exactly in $\Q(\sqrt2)$.

There is a useful independent magnitude check. Let
$B(u,v)=Q(u+v)+Q(u)+Q(v)+Q(0)$ have rank $\rho$ on $\F_2^d$, $d=2m$.
Changing variables from a pair of summation points to their difference gives
\[
 S(Q)^2=2^d\sum_{h\in\ker B}(-1)^{Q(h)+Q(0)}.
\]
The exponent restricts to a linear functional on $\ker B$. If it is
nonzero the sum vanishes; otherwise $|S(Q)|=2^{d-\rho/2}$. The elimination
algorithm supplies the sign that this square identity does not determine.
\hfill$\square$

The obstruction is not simply the absence of strongly correlated examples.
For even $m$, put
$q(x)=x_1x_2+\cdots+x_{m-1}x_m$. Summing separately over pairs shows
$\widehat f(y)=N^{-1/2}f(y)$. Hence
\[
 \Phi(f,f)=1,\qquad
 \Phi\bigl(f,\,y\mapsto f(y)(-1)^{y_1}\bigr)=0.
\]
These have maximally dense Fourier support and a large decision gap, but
Lemma~35.2 evaluates both efficiently. Thus merely combining an explicit
bent phase with an exact promise does not establish the desired hardness.

\medskip
\textbf{Where succinctness and totality intervene.}
Supplying full tables has input length $2N$. The fast Walsh transform
computes \eqref{r35:phi} in $O(N\log N)$ arithmetic operations on integers
with polynomially bounded bit lengths, so the query lower bound does not
become a superpolynomial bound in that encoding. Supplying small circuits
instead avoids the table-size cost, but allows classical algorithms to
exploit the circuit syntax, as in Lemma~35.2.

For arbitrary circuit descriptions, the test $\Phi\ge a$ versus $\Phi\le b$
with fixed $a>b$ is a promise problem. Defining the remaining inputs to be
no-instances does not make the quantum tester a total bounded-error
algorithm: amplitudes can lie arbitrarily close to the new yes/no boundary.
A successful construction must specify a total language and show its
algorithm has error at most $1/3$ on \emph{every} input, or provide a
polynomial-time recognizable encoding whose valid descriptions have a
proved gap and whose invalid descriptions receive deterministic answers.
The oracle/promise distinction emphasized in \eref{35}{3} is therefore a
live issue here, not a cosmetic qualification.

The strongest plausible continuation is an explicit nonlinear family for
which the gap is certified from syntax but no polynomial-hierarchy
algorithm decides the associated label. To establish the target it would
have to produce \emph{one} language outside every fixed level, not a
sequence $L_k$ with no uniform method for assembling them. No such family
or lower-bound argument has been found in this attempt.

\paragraph{Stress test.}
The elimination identity is checked in the accompanying script for all
quadratic polynomials in four binary variables, against exhaustive summation.
The dense examples are checked at even $m=2,4,6$. These are exact integer
checks, not tests of a complexity-class separation. The exponent-sum
algorithm is restricted to degree two; claiming its extension to arbitrary
polynomial-size Boolean circuits would discard the essential invariant
that $C+AB$ is still quadratic.

An oracle may contain information with no efficient uniform description.
Replacing it by a random seed without bounding both seed length and
evaluation cost is not de-oracularization. Fixing a successful seed
nonconstructively may introduce advice. Neither of those steps is used
here. The desired conclusion would in particular exclude
$\mathsf{BQP}\subseteq\mathsf{NP}$, since $\mathsf{NP}\subseteq\mathsf{PH}$.
No such unrelativized exclusion follows from either of the lemmas.

\Needspace{8\baselineskip}
\paragraph{Outcome.}
\textbf{Outcome: Exploratory approach with unresolved gap.}

The attempt proves two explicit obstructions to simple replacements of the
Forrelation oracle: Fourier sparsity rules out a constant signal, whereas
quadratic phases allow a constant signal but admit exact polynomial-time
classical evaluation. The remaining tasks are a uniform total-gap
construction and a lower bound against all fixed hierarchy levels that
uses more than oracle indistinguishability. Neither is supplied. The
auxiliary calculations are not asserted to be historically new, and no
complete separation candidate is claimed.

% END RESEARCH 35

\subsection{Sources}
\begin{itemize}\raggedright
\item[\textnormal{[S1]}]\hypertarget{source:35:1}{} Scott Aaronson, BQP and the Polynomial Hierarchy, arXiv:0910.4698, Abstract (2009).
\url{https://arxiv.org/abs/0910.4698}.
{\small\textit{Catalogue formulation source.}}
\item[\textnormal{[S2]}]\hypertarget{source:35:2}{} BQP and the Polynomial Hierarchy.
\url{https://www.scottaaronson.com/papers/stoc114-aaronson.pdf}.
{\small\textit{Catalogue research/status reference; not a proof endorsement.}}
\item[\textnormal{[S3]}]\hypertarget{source:35:3}{} IQP circuits for 2-Forrelation.
\url{https://arxiv.org/abs/2604.15248}.
{\small\textit{Catalogue research/status reference; not a proof endorsement.}}
\item[\textnormal{[S4]}]\hypertarget{source:35:4}{} Promises should be taken seriously: On relativization with promise problems.
\url{https://arxiv.org/abs/2609.07945}.
{\small\textit{Catalogue research/status reference; not a proof endorsement.}}
% BEGIN ADDED REFERENCES 35
\item[\textnormal{[E1]}]\hypertarget{extra:35:1}{} Ran Raz and Avishay Tal,
\emph{Oracle Separation of BQP and PH}, ECCC TR18-107 (2018);
Journal of the ACM 69(4) (2022), article 30.
\url{https://eccc.weizmann.ac.il/report/2018/107/}.
\url{https://doi.org/10.1145/3530258}.
{\small\textit{Oracle scope and distributional statement checked
27 September 2026.}}
\item[\textnormal{[E2]}]\hypertarget{extra:35:2}{} Quentin Buzet and
Andr\'e Chailloux, \emph{IQP circuits for 2-Forrelation},
arXiv:2604.15248v2, 24 August 2026, abstract and resource statement.
\url{https://arxiv.org/abs/2604.15248v2}.
{\small\textit{Version-specific supplement to [S3]; not an unrelativized
separation. Checked 27 September 2026.}}
\item[\textnormal{[E3]}]\hypertarget{extra:35:3}{} David Miloschewsky,
Supartha Podder and Dorian Rudolph, \emph{Promises should be taken
seriously: On relativization with promise problems}, arXiv:2609.07945v1,
7 September 2026, abstract and language/promise separation statement.
\url{https://arxiv.org/abs/2609.07945v1}.
{\small\textit{Recent manuscript, statement-level consultation
27 September 2026; not a complete proof audit.}}

% END ADDED REFERENCES 35
\end{itemize}

% END PRESERVED SECTION
\end{document}
